%% ============================================================
%% PAPER_B_MANUSCRIPT.tex — generated July 10, 2026
%% Multi-Scale Decoupling of Jeffrey Updating on Correlated Soft Cues
%% Target: Journal of Mathematical Economics (sole target).
%%
%% GENERATION NOTES (delegated decisions — flag for Anurag review):
%%  (a) §RL numbered as Section 3; downstream sections renumbered 1–7.
%%  (b) bibliography.bib extended with Field1978, Epstein2006,
%%      Ortoleva2012 (verification files PRESENT in bundle — hard rule
%%      satisfied). Jeffrey1983 cited from the existing shared .bib;
%%      its bundle verification file remains the top pre-submission
%%      gap per VERIFICATION_STATUS_LEDGER.md. ALL other sources
%%      appear as visible \pcite{...} pending slots, NOT in the .bib.
%%  (c) Vignette guardrail sentence (S7 content) placed in §1
%%      immediately after the vignette.
%%  (d) Benchmark naming policy: proper name "Bayes-factor benchmark"
%%      at introduction (§2.2); "matched Bayesian posterior" only
%%      where coincidence with Bayesian conditionalization is itself
%%      the content (Prop. 1); "the benchmark" elsewhere.
%%  Proof placement per PROOF_PLACEMENT_MAP.md (agreed July 10):
%%  main text A1/B1/C1/C2 in full; A2 statement + mechanism paragraph;
%%  B2 statement + band sketch; full proofs of A2, B2, C3 in
%%  Appendix A. This supersedes the skeleton's earlier "no proof
%%  appendix" line. Density-robustness and sign-robustness remarks
%%  stay in the main text after Theorem 1 (objections-as-remarks
%%  discipline).
%%
%% elsarticle.cls is not installed in this runtime; compile locally.
%%
%%  STRUCTURAL PASS (bottom-up tightening, applied this session):
%%   - Four assumptions hoisted from Scope to a new Setup subsection
%%     (\S Assumptions), stated bare; Scope kept as a backward-looking
%%     robustness recap (what each assumption is load-bearing for).
%%   - Prop. C3 promoted to its own subsection ("The protected class"),
%%     the visible capstone, immediately before Incidence-vs-intensity.
%%   - Setup->A1 and rem:density->C2 refs re-signposted as explicit
%%     forward previews ("established below", "below").
%%   - Related literature kept at \S3 (econ convention, author's call).
%% ============================================================

\documentclass[review,12pt,authoryear]{elsarticle}

\usepackage{amsmath,amssymb,amsthm}
%% lineno removed: not a JME (Journal of Mathematical Economics) requirement, per
%% https://www.sciencedirect.com/journal/journal-of-mathematical-economics/publish/guide-for-authors
%% (checked July 10, 2026). It was inherited from MANUSCRIPT_TEMPLATE.tex without audit against
%% the actual target-journal guide; likely confused with the Journal of MONETARY Economics,
%% which is the journal that does require line numbers.

%% ---- theorem environments (separate counters; skeleton numbering:
%%      Propositions 1–5, Lemma 1, Theorem 1) ----
\newtheorem{proposition}{Proposition}
\newtheorem{lemma}{Lemma}
\newtheorem{theorem}{Theorem}
\theoremstyle{remark}
\newtheorem{remark}{Remark}
\theoremstyle{definition}
\newtheorem{assumption}{Assumption}

%% ---- notation macros (mirror proofs_selfcontained.tex) ----
\newcommand{\PB}{P^{\mathrm{B}}}
\newcommand{\PJ}{P^{\mathrm{J}}}
\newcommand{\Pbar}{\overline{P}}
\newcommand{\assoc}{\operatorname{assoc}}
\newcommand{\sgn}{\operatorname{sgn}}
\newcommand{\E}{\mathbb{E}}
\newcommand{\bigO}{\mathcal{O}}
\newcommand{\vv}{v}
%% \anchor is a source-only grounding marker. It preserves the argument in the
%% .tex file so the grounding audit (GROUNDING_MAINTENANCE_PLAN.md) can grep for
%% \anchor{TAG} — but it expands to nothing in the compiled PDF, so no external
%% "Proofs" document is cited to the reader. The manuscript is self-contained.
\newcommand{\anchor}[1]{}

%% ---- pending-citation slot: visible, greppable; NOT a .bib entry.
%%      Every \pcite must be resolved (verification file written,
%%      entry added to bibliography.bib, \pcite -> \citep/\citet)
%%      before submission. See VERIFICATION_STATUS_LEDGER.md. ----
\newcommand{\pcite}[1]{\textbf{[#1 --- pending verification]}}

\journal{Journal of Mathematical Economics}

\begin{document}

\begin{frontmatter}

\title{Multi-scale decoupling of Jeffrey updating on correlated soft cues}

\author[rvt]{Anurag Srivastava\corref{cor1}}
\ead{[author email to insert]}

\cortext[cor1]{Corresponding author.}

\affiliation[rvt]{organization={Riskcare Ltd.},
                  addressline={[affiliation address to insert]},
                  city={London},
                  country={United Kingdom}}

\begin{abstract}
We study the aggregation of Jeffrey updating on two correlated binary soft cues. On a single
cue, the Jeffrey posterior coincides with the matched Bayesian posterior identically; on two
cues the composite update is non-commutative at first order in the prior covariance $c$ between
the attributes. We prove a separation of scales: individual belief and decision diverge from the
benchmark at $\bigO(c)$, while the two quantities a group-level audit characteristically
computes---the mean believed association and the surplus-weighted welfare loss---are
$\bigO(c^{2})$. The protected class is characterised exactly: a smooth statistic of the mean
posterior is $\bigO(c^{2})$-close to the benchmark across an open set of priors if and only if
its differential at independence is a multiple of $d\,\assoc$; the protected class comprises the
local dependence measures and, for the threshold decision, the surplus-weighted welfare loss.
The headcount of sequence-affected individuals remains $\bigO(c)$, so an aggregate matching the
benchmark does not license the inference that its members are Bayesian.
\end{abstract}

%% Highlights removed per author (2026-07-11): JME's own guide could not be fetched to
%% confirm whether Highlights are requested; block commented out rather than deleted so it
%% is recoverable if the guide turns out to require them. If restored, bullet 2 is already
%% trimmed to <=85 chars, and Highlights must ship as a separate editable file at submission.
%\begin{highlights}
%\item Sequential Jeffrey updating on correlated soft cues is $O(c)$ order-dependent.
%\item Mean association and surplus-weighted welfare loss are $O(c^2)$; headcount is $O(c)$.
%\item The believed association is the unique smooth statistic protected to second order.
%\item Protection is structural (separability), not statistical (noise cancellation).
%\item Aggregate Bayes-conformity does not certify member-level Bayesian rationality.
%\end{highlights}

\begin{keyword}
Jeffrey conditioning \sep probability kinematics \sep belief aggregation \sep scale separation \sep order effects \sep welfare measurement
\JEL C11 \sep D81 \sep D83 \sep D63
\end{keyword}

\end{frontmatter}

\section{Introduction}
\label{sec:intro}

That the order in which evidence arrives can move the final judgment is among the oldest
findings in the study of impression formation \pcite{Asch 1946}, and order effects in sequential
belief revision are documented across tasks and elicitation modes \pcite{Hogarth--Einhorn 1992}.
This paper does not add to that documentation. It takes one coherent mechanism for
order-dependence and traces it from the individual to the population, and one direction of the
tracing deserves emphasis at the outset: the mechanism studied here \emph{suffices} for
order-dependence; observed order-dependence does not certify the mechanism. Processes with
overlapping signatures exist---anchoring-and-adjustment models of belief revision produce order
effects with no Jeffrey machinery \pcite{Hogarth--Einhorn 1992}---so nothing in what follows
diagnoses any observed order effect as the mechanism modelled here. The results establish
sufficiency, never diagnosis.

Consider a hiring panel working through a candidate's file. Two pieces of evidence bear on the
decision: a credential that speaks to competence, and the tone of a reference letter that speaks
to trustworthiness. Neither is decisive; each leaves an impression. One panel happens to read
the credential first and the letter second; an otherwise identical panel reads them in the
reverse order. If the two attributes are believed to go together at all, the two panels finish
with different beliefs about the same candidate---and, if the candidate sits near the panel's
threshold, with different decisions. The difference is driven by nothing in the evidence and
everything in its sequencing. Two clarifications belong next to the example rather than at the
end of the paper. First, the effect modelled here is idiosyncratic to the encounter: it reverses
if the reading order flips, and no attribute of the candidate predicts its direction, so it is
not the systematic, marker-keyed phenomenon that the term ``discrimination'' ordinarily denotes.
Second, the example exhibits a configuration; the paper asserts nothing about how often the
configuration is realised in actual evaluation settings.

The stake is nonetheless real for the individual concerned: identical evidence can license
different decisions, with the difference set by arrival order rather than by anything the
evidence contains.

What the vignette identifies is an \emph{input type}, not an updating rule. The panelist's
impressions arrive as revised credences, not as likelihoods: the panelist comes to feel the
candidate competent, not to record that this credential is some calibrated multiple more likely
under competence than under its absence. For input of that kind, strict conditionalization is
not merely inferior but unavailable---there is no certifiable proposition to condition on. And
given a delivered credence on a cue's partition, with the experience bearing on nothing beyond
that partition, Jeffrey conditioning is the unique coherent revision: holding the conditionals
fixed is equivalent to making the minimal change of the prior consistent with the delivered
marginal \citep{DiaconisZabell1982}. The rule itself---reset the cue's credence, hold the
conditionals fixed---is \citet[ch.~11]{Jeffrey1983}; the theorem-level grounding is collected in
Section~\ref{sec:jeffrey}. Each single Jeffrey step, moreover, coincides exactly with a matched
Bayesian conditionalization (Proposition~\ref{prop:A1}), so no individual step deviates from
Bayes. Yet on two cues in sequence the composite update is non-commutative whenever the two
attributes are believed to go together: reversing the order shifts the posterior, and any
threshold decision it drives, in proportion to the prior covariance $c$ between the
attributes---the strength with which the two traits are believed to go together before any cue
arrives (Proposition~\ref{prop:A2}). That successive Jeffrey updates are asymmetric in the
delivered-credence parametrisation is \citet{Field1978}'s observation; that they commute when
their inputs are held fixed in Bayes-factor form is \citet{Wagner2002}'s theorem. The order
effect is therefore a property of the input type, not a defect of the rule.

The aggregate question inherits its premise from the input type. Much of institutional
evaluation delivers credences rather than likelihoods, so a population of evaluators is a
population of Jeffrey updaters by the nature of its evidence, not by modelling choice; encounter
order is idiosyncratic to each evaluator, so the population realises a distribution over orders;
and what is observed institutionally is rarely an individual posterior but an aggregate. Which
aggregates carry the individual effect is therefore the operative question.

The question of the paper is what this individual sequence-dependence does as it passes to the
population. The answer is that it does not vanish; it changes scale, and it does so selectively.
Every individual diverges from the benchmark at first order in $c$
(Proposition~\ref{prop:A2}). Yet the two quantities a group-level audit characteristically
computes---the mean believed association between the attributes
(Proposition~\ref{prop:C1}) and the surplus-weighted welfare loss of the population's decisions
(Theorem~\ref{thm:B2})---are second order. The suppression of the first-order effect is not a
pair of coincidences but a characterization (Proposition~\ref{prop:C3}): the class of statistics
protected to second order is exactly delimited. It contains the local dependence measures---the
association and, to first order at independence, any smooth function of it---together with the
surplus-weighted welfare loss. Nothing else does: every marginal and every nonconstant-weight
mean is genuinely first order, and the headcount of sequence-affected individuals is first order
as well (Proposition~\ref{prop:C2}). The instrumental consequence is immediate: an instrument
built from the protected aggregates has one order less resolution than the phenomenon it would
detect, and what an audit registers is settled by which statistic it aggregates, not by how much
data it collects. Put in terms of rationality rather than measurement: a population can look
aggregate-Bayesian---benchmark-conforming on every protected statistic---while no member is
Bayesian, each diverging at first order.

The scope is deliberately minimal. The model has two binary attributes, two attribute-local soft
cues, and works at leading order in the feasibility-bounded prior covariance $c$; two cues is
the smallest configuration in which order can matter at all, and extension beyond it is posed as
an open direction in Section~\ref{sec:conclusion}. The model exhibits the configuration and
asserts nothing about its real prevalence. Section~\ref{sec:literature} situates the result
against the belief-aggregation, sequential social-learning, non-Bayesian-updating, and
statistical-discrimination literatures.

The paper proceeds as follows. Section~\ref{sec:setup} sets up the model and the benchmark.
Section~\ref{sec:literature} relates the paper to the literature.
Section~\ref{sec:individual} establishes the individual-level results: exact immunity on a
single cue and first-order divergence on two. Section~\ref{sec:aggregation} contains the
separation of scales: the welfare law, the decoupling of the believed association, the
characterization of the protected class, and the criterion split.
Section~\ref{sec:scope} states the domain and what would break the results, and
Section~\ref{sec:conclusion} concludes. Proofs of Propositions~\ref{prop:A2}
and~\ref{prop:C3} and Theorem~\ref{thm:B2} are in \ref{app:proofs}; the remaining proofs are
short and given in the text.

\section{Setup}
\label{sec:setup}

\subsection{Decision problem and welfare criterion}
\label{sec:decision}

Two binary attributes $(A,B)\in\{0,1\}^{2}$---competence and trustworthiness, say---are each
present or absent, so a belief assigns probability to the four configurations at once: it is a
\emph{joint law}, a single distribution over the pairs $(i,j)$---a point in the probability
simplex $\Delta^{3}$---rather than a pair of one-attribute distributions \anchor{Setup}. The
distinction motivates the term and is the load-bearing choice. The two marginal laws record only
how likely each attribute is on its own, and fixed marginals leave a whole family of joint laws
open, its members differing precisely in how the attributes are believed to co-move; a
cross-attribute association therefore exists only on the joint law, never on its two marginals
separately.
The prior $P$ is the evaluator's subjective credence, before any cue arrives, over the joint
attributes of a randomly drawn candidate. We coordinatize $\Delta^{3}$ by the marginals
$P(A{=}0)=\alpha$, $P(B{=}0)=\beta$ and the covariance $c$---a change of coordinates on the
simplex, not an added assumption:
\[
  P=\begin{pmatrix}\alpha\beta+c & \alpha(1-\beta)-c\\[2pt](1-\alpha)\beta-c & (1-\alpha)(1-\beta)+c\end{pmatrix},
  \qquad \operatorname{Cov}(A,B)=c,
\]
so $c=0$ is independence and $c$ is feasibility-bounded by nonnegativity of the four cells
\anchor{Setup}. The parameter $c$ is the prior covariance between the attributes; it is also the
prior value of the believed association defined in Section~\ref{sec:association}---the same
quantity read before and after updating.

The evaluator receives two \emph{soft cues}, and the cues are attribute-local: a cue on $A$
yields an impression (a target marginal) $q=(q_{0},q_{1})$ on $A$'s partition, a cue on $B$
yields an impression $r=(r_{0},r_{1})$ on $B$'s partition, and no cue bears on a joint event.
Each impression is a delivered credence, not a likelihood. The independent product of the two
impressions is written $q\otimes r$, with $(q\otimes r)_{ij}=q_{i}r_{j}$ \anchor{Setup}.

A desirability vector $\vv=(v_{ij})_{i,j\in\{0,1\}}$---one entry per state-space cell, the same shape as $P$---gives the score
$s(P)=\langle\vv,P\rangle=\sum_{ij}v_{ij}P(i,j)$; the evaluator engages if and only if
$s(P)\ge\tau$ \anchor{Setup}. Thresholding the score is expected-utility maximization at a given
belief: $s(P)$ is the expected desirability of engaging and $\tau$ the value of the outside
option. We take the posterior $\PB$---the Bayes-factor benchmark, constructed in
Section~\ref{sec:jeffrey}---as the welfare benchmark. Its status is declared, not asserted: it conditions once on the impressions'
likelihood content matched from the prior marginals \citep{Wagner2002}, it coincides with the
single-impression Jeffrey update (established below as Proposition~\ref{prop:A1}), and it is order-free by
construction; the \emph{welfare-optimal} action is \emph{defined} as engaging exactly when
$s(\PB)\ge\tau$, inherits these three properties, and asserts no ground-truth optimality beyond
them \anchor{Setup}. Write $u:=s(\PB)-\tau$, the net surplus of engagement relative to the
benchmark; the cost of the action that differs from the benchmark action, so measured, is
$|u|$.

The welfare object of the paper is the population aggregate of these stakes:
$L(c)$, the expected foregone surplus, across evaluators and the sequence distribution, of the
Jeffrey decision relative to the benchmark decision \anchor{Thm.~B2}---that is, the expectation
of $|u|$ over those whose decision the sequence actually changes, and zero for everyone else. As
the band picture following Theorem~\ref{thm:B2} shows, that restriction is what turns an
$\bigO(c)$ individual effect into an $\bigO(c^{2})$ aggregate. It is the evaluators' own
criterion (expected desirability) evaluated at the declared benchmark law---a single-agent
regret quantity, not a Pareto, total-surplus, or social-welfare-function object: no allocation
is traded across agents and no first-best efficiency is claimed. The weight $|u|$ is a surplus
weight, not a penalty convention, and the choice of criterion is load-bearing: it is exactly
what Section~\ref{sec:criterion} shows to control the leading order of the aggregate.

\subsection{Jeffrey conditioning and the Bayes-factor benchmark}
\label{sec:jeffrey}

Jeffrey's rule updates on an impression by resetting the cue's marginal to the delivered
credence and holding the conditionals fixed---the invariance condition
\citep[ch.~11]{Jeffrey1983}. For an impression $q$ on $A$,
\[
  (\PJ_{A}P)(i,j)=q_{i}\,P(i,j)/P(A{=}i),
\]
and symmetrically for an impression $r$ on $B$ \anchor{Setup}. The axiomatic grounding is
twofold. The partition satisfying the invariance condition is the minimal sufficient statistic
for revising the prior to any candidate posterior on that partition
\citep[Thm.~2.2]{DiaconisZabell1982}, and the rule is the minimal revision of the prior
consistent with the delivered marginal, in the $I$-divergence sense
\citep[Thm.~5.1]{DiaconisZabell1982}: the update that changes as little else as possible. The
two-cue Jeffrey posteriors are the two sequences
\[
  \PJ_{AB}=\PJ_{B}\!\circ\PJ_{A}\,P,\qquad \PJ_{BA}=\PJ_{A}\!\circ\PJ_{B}\,P ,
\]
and the order effect, the object of study at the individual level, is
$\PJ_{AB}-\PJ_{BA}$ \anchor{Setup}.

The benchmark $\PB$ is not a rival updating rule but the same cues processed in a different
currency. Each impression's evidential content can be written as a Bayes factor---a likelihood
matched to the cue given the prior, $\ell^{A}_{i}\propto q_{i}/P(A{=}i)$ and
$\ell^{B}_{j}\propto r_{j}/P(B{=}j)$---and, unlike a delivered credence, a Bayes factor carries
the same content whatever prior it meets. Combining the two cues' Bayes-factor content is
order-invariant \citep{Wagner2002}, and $\PB$ is the single update on that combined content,
\[
  \PB(i,j)\propto P(i,j)\,\ell^{A}_{i}\,\ell^{B}_{j},
\]
well-defined regardless of which cue arrived first \anchor{Setup}; on any one cue it coincides
with the matched Jeffrey step (Proposition~\ref{prop:A1}). The comparison of the paper is thus
internal to probability kinematics: sequential updating on delivered credences, which does not
commute, against one-shot updating on Bayes-factor content, which does.

\subsection{Assumptions}
\label{sec:assumptions}

The results hold on a deliberately narrow domain, fixed by the four assumptions below;
Section~\ref{sec:scope} states what each is load-bearing for and what would break if it were
dropped.

\begin{assumption}[Locality in $c$]\label{as:localc}
Two binary correlated cues; the claims are at leading order in the prior covariance $c$ between
the attributes, which is feasibility-bounded by cell nonnegativity. The separation of scales is
a local statement in $c$.
\end{assumption}

\begin{assumption}[Soft cues]\label{as:soft}
The cues are bundled, correlated, and soft: each delivers a credence on an attribute's partition,
not a decisive marker.
\end{assumption}

\begin{assumption}[Attribute-locality]\label{as:local}
Each impression is a credence on a single attribute's partition, never on a joint event.
\end{assumption}

\begin{assumption}[Surplus criterion]\label{as:surplus}
Welfare is measured by foregone surplus (cost $|u|$), with a threshold density that is at least
integrable.
\end{assumption}

\section{Related literature}
\label{sec:literature}

Four literatures border the result; in each case the surface is shared and the object is not.

\emph{Belief aggregation.} A normative literature evaluates group belief formation by whether
aggregation commutes with updating---a design criterion for the pooling rule
\pcite{Dietrich 2021}. Geometric pooling satisfies that criterion; the population mean studied
here is linear pooling and fails it generically, consistent with the first-order divergence of
every statistic outside the protected class of Proposition~\ref{prop:C3}. The present paper is
positive rather than normative: it takes the mean belief as the object an observer computes and
characterises which of its statistics conceal the members' order-dependence.

\emph{Sequential social learning.} Economics has previously admitted order-dependence at the
population level through observational learning: herds and informational cascades
\pcite{Banerjee 1992; Bikhchandani--Hirshleifer--Welch 1992}. The mechanism and the object both
differ from the present ones, on four axes. The sequencing there is across agents (the order in
which predecessors act), not within an agent (the order in which one evaluator meets two cues).
The coupling is a feed-forward informational externality---agent $n$'s action enters agent
$n{+}1$'s information set---whereas the evaluators here are mutually independent, with no
feedback. The order-dependence there is path-dependence of a temporal accumulation: which
signals happen to arrive early get frozen in, because a coarse public action fails to be a
sufficient statistic for private information---when the action is rich enough to invert, order
washes out \pcite{Banerjee 1992}; a cascade is precisely an action that conveys no private
signal \pcite{Bikhchandani--Hirshleifer--Welch 1992}. Here, by contrast, the object is literal
permutation non-commutativity of a fixed pair of inputs read by one agent
(Proposition~\ref{prop:A2}). And the state there is single-valued, so no joint object exists:
their order effect lives on the analogue of the marginals, which are unprotected here as well;
the believed association, the dependence object this paper's protection theorem concerns, has no
counterpart in single-state models.

\emph{Non-Bayesian updating.} An axiomatic literature makes departure from Bayes itself the
primitive: \citet{Epstein2006} makes the updating rule subjective, and \citet{Ortoleva2012}
axiomatizes departures triggered by unexpected news. Here the rule is held fixed---Jeffrey's,
the coherent rule for likelihood-free input---and the object is its aggregation signature. Work
on the commutativity of sequential Jeffrey updating restores order-invariance by changing how
successive inputs are pooled \pcite{Pettigrew--Weisberg 2025}; this paper instead accepts the
non-commutativity and characterises what survives of it in aggregate. Both exercises are
individual-level in their subject matter; the population axis is the present contribution.

\emph{Statistical discrimination.} The paper borrows the evaluator-with-binary-attributes frame
from the statistical-discrimination lineage \pcite{Phelps 1972; Arrow 1973}, but its object is
kinematic, not an equilibrium fixed point: in \pcite{Coate--Loury 1993} the stereotype is a
self-confirming steady state of best responses, while here it is a statistic of a belief
trajectory under a fixed updating rule. It is likewise distinct from stereotypes as
representativeness-distortion or selective recall \pcite{Bordalo--Coffman--Gennaioli--Shleifer
2016}: the believed association studied here is produced by coherent updating on impoverished
input, not by a distortion of memory or sampling. The lineage distinction the concluding section
trades on---a statistical mechanism that suppresses an aggregate signature, rather than a
structural one---goes back to \pcite{Becker 1962}. Finally, the incidence/intensity vocabulary
of Section~\ref{sec:criterion} borrows a distinction from the poverty-measurement literature
\pcite{Foster--Greer--Thorbecke 1984}---the distinction only, not the axiomatic apparatus.

\section{Order effects at the individual level}
\label{sec:individual}

\begin{proposition}[A1: exact immunity]\label{prop:A1}
On a single cue, the Jeffrey posterior equals the matched Bayesian posterior identically, for
every prior and every impression. With two cues but $c=0$ (independent attributes),
$\PJ_{AB}=\PJ_{BA}=\PB$ exactly. \anchor{Prop.~A1}
\end{proposition}

\begin{proof}
Single cue on $A$: the matched likelihood $\ell^A_i=q_i Z/P(A{=}i)$, with $Z$ the normaliser
fixing the posterior $A$-marginal to $q$, gives
$\PB(i,j)=P(i,j)\ell^A_i/Z=q_i\,P(i,j)/P(A{=}i)=(\PJ_A P)(i,j)$. For $c=0$ the prior factorises,
$P(i,j)=P(A{=}i)P(B{=}j)$; each Jeffrey step and the Bayesian update then act on independent
coordinates, so both sequences and the Bayesian posterior equal the product $q\otimes r$.
\end{proof}

Thus a decision turning on \emph{one} decisive marker is immune; the phenomenon requires a
bundle of correlated soft cues, and the single-marker case lies outside the claim's domain (Assumption~\ref{as:soft}).

\begin{proposition}[A2: micro divergence]\label{prop:A2}
For $c\neq0$ the individual gap and the sequence effect are first order and vanish at $c=0$:
\[
  \PJ_{AB}-\PB=c\,D+\bigO(c^{2}),\quad D\not\equiv 0;\qquad
  \PJ_{AB}-\PJ_{BA}=c\,\Delta_{\mathrm{seq}}+\bigO(c^{2}),\quad \Delta_{\mathrm{seq}}\not\equiv 0,
\]
with $\PJ_{AB}-\PJ_{BA}\big|_{c=0}=0$ (no constant term). Writing
$Z:=\alpha\beta(1-\alpha)(1-\beta)>0$, the leading individual gap has the rank-one form
\[
  D=\partial_c\!\left(\PJ_{AB}-\PB\right)\big|_{c=0}
   =\kappa\; q\otimes(1,-1),
  \qquad
  \kappa=\frac{(\alpha-q_0)\,r_0(1-r_0)}{Z}.
\]
\anchor{Prop.~A2}
\end{proposition}

The mechanism behind Proposition~\ref{prop:A2} is visible without the expansion that proves it.
A Jeffrey update on the $A$-cue rescales the joint distribution row by row, preserving the
conditional distribution of $B$ given $A$ \anchor{Setup}. When the attributes are independent,
preserving that conditional means leaving $B$'s marginal untouched: the two orders trivially
commute and coincide with the benchmark exactly (Proposition~\ref{prop:A1}). When
$c\neq0$, rescaling rows changes $B$'s marginal---the update reaches the second attribute
through the covariance---and the two orders no longer agree, at first order in $c$
\anchor{Prop.~A2}. The scalar $\kappa$ exhibits the two ways the leading gap can vanish:
$\kappa\propto(\alpha-q_0)$, the surprise in the $A$-impression relative to the prior, and
$\kappa\propto r_0(1-r_0)$, the spread of the $B$-impression; the gap vanishes only when the
$A$-impression coincides with the prior $A$-marginal or an impression is degenerate
\anchor{Prop.~A2}. The full expansion, the explicit sequence-effect coefficient, and the
non-vanishing certificate are in \ref{app:proofs}.

\begin{remark}[structure of the leading terms]\label{rem:structure}
The leading individual gap preserves the $A$-marginal ($D_{i0}=-D_{i1}$ for each $i$): to first
order $\PJ_{AB}$ and $\PB$ differ only in the $B$-conditional within each $A$-level. The order
effect is the genuinely two-dimensional residue. On the symmetry locus
$\{\alpha=\beta,\ q_0=r_0\}$ it survives in the antisymmetric off-diagonal part,
$\Delta_{\mathrm{seq},01}=-\Delta_{\mathrm{seq},10}
=q_0(1-q_0)(q_0-\alpha)/(\alpha^{2}(1-\alpha)^{2})$---moving scores whenever
$v_{01}\neq v_{10}$---and vanishes entirely only at the uninformative point $q_0=\alpha$.
``First impressions stick'' is the statement $\Delta_{\mathrm{seq}}\neq0$: once the attributes
covary, the earlier cue's influence is not erased by the later one, and the residue is
$\bigO(c)$. \anchor{Prop.~A2}
\end{remark}

\begin{remark}[economic reading at the individual level]\label{rem:econmicro}
Read Proposition~\ref{prop:A2} through the hiring vignette of the introduction: two otherwise
identical panels work the same file---a credential bearing on competence, a reference letter
bearing on trustworthiness---in opposite orders. If the panels take the attributes to be
independent, the two orders agree with each other and with the benchmark exactly
(Proposition~\ref{prop:A1}); everything below is proportional to $c$, the strength with which
competence and trustworthiness are believed to go together before the file is opened
\anchor{Prop.~A2}. The scalar $\kappa$ then prices the configuration. The panels depart from
the benchmark at first order only when the credential genuinely surprises---$q_{0}\neq\alpha$,
the impression departing from the prior base rate---and only when the letter leaves a genuinely
soft impression---$r_{0}(1-r_{0})>0$, neither glowing nor damning enough to be decisive
\anchor{Prop.~A2}; a credential that merely confirms what the panel already assumed, or a
letter read as dispositive, kills the leading term. Where the departure lives is equally
specific (Remark~\ref{rem:structure}): to first order the credential-first panel holds the
benchmark view of \emph{how likely} competence is and departs only in what competence is taken
to \emph{imply about} trustworthiness---the conditional reading of the file, not its headline
rate. The sequence effect is the two panels' disagreement with each other: whichever evidence a
panel happens to read first leaves a residue the later evidence does not erase, and even in the
fully symmetric configuration ($\alpha=\beta$, $q_{0}=r_{0}$) the residue survives in the
antisymmetric cells, moving the score whenever $v_{01}\neq v_{10}$
(Remark~\ref{rem:structure}). Its direction reverses when the reading order flips
\anchor{Prop.~A2}---set by the encounter and the cue--prior geometry, not by any attribute of
the candidate, as the introduction's clarification already fixed. What that moving score costs
when it crosses a threshold is the subject of Section~\ref{sec:aggregation}.
\end{remark}


\section{Aggregation and the separation of scales}
\label{sec:aggregation}

Every individual leaves Section~\ref{sec:individual} at first order from the benchmark
(Proposition~\ref{prop:A2}); this section asks what of that divergence a population shows. The
naive expectation is inheritance: the score perturbation entering each threshold decision is
linear in the individual gap (Lemma~\ref{lem:B1} below), so a first-order cause should leave a
first-order aggregate trace. The results of the section are that it does not, and that the
suppression is selective. The surplus-weighted welfare loss is second order
(Theorem~\ref{thm:B2}), and so is the mean believed association (Proposition~\ref{prop:C1});
the protection is then shown to be a characterization rather than a pair of coincidences---to
first order, the association is the only statistic aggregation protects
(Proposition~\ref{prop:C3})---and the boundary is drawn by the headcount of sequence-affected
individuals, which remains first order (Proposition~\ref{prop:C2}). The ascent runs from
welfare through belief to the characterization, and closes on the split the characterization
explains.

\subsection{Welfare}
\label{sec:welfare}

The score perturbation entering welfare is inherited from Proposition~\ref{prop:A2}, not
posited.

\begin{lemma}[B1: inherited perturbation]\label{lem:B1}
$s(\PJ_\sigma)-s(\PB)=c\,\delta_\sigma+\bigO(c^{2})$ with
$\delta_\sigma=\big\langle\vv,\ \partial_c(\PJ_\sigma-\PB)|_{c=0}\big\rangle$; generically
$\delta_{AB}\neq\delta_{BA}$, and $\delta_\sigma=0$ if $\vv$ is constant. \anchor{Lem.~B1}
\end{lemma}

\begin{proof}
Apply the linear functional $s=\langle\vv,\cdot\rangle$ to
$\PJ_\sigma-\PB=cD_\sigma+\bigO(c^2)$ from Proposition~\ref{prop:A2}. Constancy of $\vv$ makes
$\delta_\sigma$ telescope to $0$ since both posteriors are normalised.
\end{proof}

\begin{theorem}[B2: welfare law]\label{thm:B2}
Let $u=s(\PB)-\tau$ have, across the population and the sequence distribution, a joint density
whose $u$-marginal is integrable (Assumption~\ref{as:surplus}). Then the expected welfare loss $L(c)$ of the Jeffrey
population relative to the benchmark is $o(c)$---strictly higher order than the individual
divergence. If moreover the $u$-marginal is bounded and continuous at $u=0$ with value
$f(0)>0$, the loss is exactly second order, with constant named:
\[
  L(c)=\frac{f(0)}{2}\,\E\!\big[\delta_\sigma^{2}\mid u=0\big]\,c^{2}+\bigO(c^{3}).
\]
\anchor{Thm.~B2}
\end{theorem}

The proof, given in full in \ref{app:proofs}, turns on a band picture. The Jeffrey and benchmark
actions differ only when the perturbation $c\delta_\sigma$ flips the sign of the surplus $u$, so
the flip set is confined to a band of width $\bigO(c)$ at the threshold; on that band the cost
$|u|$ is itself $\bigO(c)$. An $\bigO(c)$ mass of individuals, each costing $\bigO(c)$, yields an
$\bigO(c^{2})$ loss; integrability alone already forces $L(c)=o(c)$, and continuity of the
density at the threshold names the constant \anchor{Thm.~B2}.

\begin{remark}[the smoothness hypothesis is a convenience, not a load-bearing assumption]\label{rem:density}
Only integrability is needed for the decoupling ($L=o(c)$, strictly above the $\bigO(c)$
headcount of Proposition~\ref{prop:C2} below); smoothness merely names the constant. The natural
failures do not break it (computationally verified):
(i) an \emph{atom exactly at the threshold} is harmless---those candidates have $|u|=0$, so
mis-deciding them costs nothing under the surplus criterion;
(ii) an integrable \emph{blow-up} $f(u)\sim|u|^{-\gamma}$, $0<\gamma<1$, degrades the rate to
$\bigO(c^{2-\gamma})$, still $o(c)$;
(iii) a \emph{vanishing} density $f(0)=0$ makes the loss $o(c^2)$, strengthening the decoupling.
The only route to a first-order welfare loss is not a density pathology but a change of
criterion (surplus to headcount, Proposition~\ref{prop:C2}): no density can route the surplus
loss to $\bigO(c)$, since a cost of $|u|=\bigO(c)$ on the band multiplies an $\bigO(c)$ band
mass. \anchor{Thm.~B2}
\end{remark}

\begin{remark}[sign-robustness]\label{rem:sign}
$\delta_\sigma$ enters \emph{squared}: the order is $c^{2}$ whether the cue sequence is random
($\E\delta_\sigma=0$) or systematic ($\E\delta_\sigma\neq0$). The sequence distribution sets
only the constant. \anchor{Thm.~B2}
\end{remark}


\subsection{Believed association}
\label{sec:association}

Let $\Pbar=\tfrac12(\PJ_{AB}+\PJ_{BA})$ be the population's mean (random-sequence) belief, and
$\assoc(D)=D_{00}D_{11}-D_{01}D_{10}$ the association of a joint law $D$; the prior value of the
association is $c$ (Section~\ref{sec:decision}). A \emph{stereotype}, in the sense of this
paper, is a believed association absent from---or in excess of---the benchmark: a nonzero
$\assoc(\Pbar)-\assoc(\PB)$ \anchor{Prop.~C1}.

The engine of the dependence-side protection is a structural fact worth stating once, at full
generality, because nothing in it is specific to Jeffrey conditioning or to two cues. It
carries three loads in what follows. It is the operative step in the proof of
Proposition~\ref{prop:C1}: every updating route in the model composes to a separable
reweighting, and a separable reweighting can only rescale an association, never shift it. It is
the source of the result's two morals---that the suppression is structural rather than a
cancellation of terms, and that it is not Jeffrey's property but attribute-locality's, holding
for any updating rule whose steps read one attribute at a time
(Remark~\ref{rem:mechanism})---and, because it holds at every $N$, of the general-$N$ reach
claimed in Section~\ref{sec:conclusion}. And it is what makes Assumption~\ref{as:local}
load-bearing rather than decorative (Section~\ref{sec:scope}). The welfare law rests on a
different mechanism entirely---the band picture of Theorem~\ref{thm:B2}---so the lemma is the
engine of the belief-level results, not of the section wholesale.

\begin{lemma}[separability of attribute-local composites]\label{lem:sep}
Fix any number $N$ of binary attributes and a joint prior $P$ over $\{0,1\}^{N}$. Any finite
sequence of attribute-local updates---each multiplying the law by a factor that reads the
configuration through a single attribute's coordinate, as every Jeffrey update on one
attribute's partition does---produces a \emph{separable} reweighting:
\[
  P'(x)=P(x)\prod_{a=1}^{N}g_{a}(x_{a})
\]
for some single-attribute functions $g_{a}$; equivalently, $\log(P'/P)$ contains no interaction
term of any order. \anchor{Lem.~SEP}
\end{lemma}

\begin{proof}
Write the update sequence as $P_{0}=P$ and, for $t=1,\dots,T$,
\[
  P_{t}(x)=P_{t-1}(x)\,h_{t}(x_{a_{t}}),
\]
where $a_{t}$ is the attribute the $t$-th update reads and $h_{t}$ is a nonnegative function
of that single coordinate. Jeffrey updates have exactly this form: the step on attribute $a$
to impression $q$ multiplies each cell mass by $h(x_{a})=q_{x_{a}}/P_{t-1}(A_{a}{=}x_{a})$,
and although this factor is state-dependent---its normalizer reads the \emph{current} law's
marginal---it depends on the configuration $x$ only through the coordinate $x_{a}$, which is
all the argument uses.

We claim, by induction on $t$, that
$P_{t}(x)=P(x)\prod_{a=1}^{N}g^{(t)}_{a}(x_{a})$ for some single-attribute functions
$g^{(t)}_{a}$. The base case $t=0$ holds with $g^{(0)}_{a}\equiv1$. For the inductive step,
suppose the claim at $t-1$; then
\[
  P_{t}(x)
  =P(x)\Big[\prod_{a}g^{(t-1)}_{a}(x_{a})\Big]h_{t}(x_{a_{t}})
  =P(x)\prod_{a}g^{(t)}_{a}(x_{a}),
  \qquad
  g^{(t)}_{a}:=
  \begin{cases}
    g^{(t-1)}_{a}\cdot h_{t}, & a=a_{t},\\[2pt]
    g^{(t-1)}_{a}, & a\neq a_{t}:
  \end{cases}
\]
the new factor folds into the $a_{t}$-th function and every other function is untouched. No
step ever couples two coordinates, so the separable form survives each step; taking
$g_{a}:=g^{(T)}_{a}$ gives the displayed form.

For the equivalent formulation, on the common support (cells with $P(x)>0$ and every
accumulated factor positive) take logarithms:
\[
  \log\frac{P'(x)}{P(x)}=\sum_{a=1}^{N}\log g_{a}(x_{a}),
\]
an additive function of the individual coordinates. In the log-linear expansion of a function
on $\{0,1\}^{N}$---the unique decomposition
$\varphi(x)=\lambda_{\emptyset}+\sum_{a}\lambda_{a}(x_{a})
+\sum_{a<b}\lambda_{ab}(x_{a},x_{b})+\cdots$ under the usual identifiability
conventions---a sum of single-coordinate functions has every term of order two and higher
identically zero, and conversely. ``Separable reweighting'' and ``no interaction term of any
order in $\log(P'/P)$'' are therefore the same condition. Nothing in the argument is specific
to Jeffrey updating, to $N=2$, or to the length of the update
sequence.\footnote{The lemma is additionally
machine-verified at general $N$ in Lean~4/Mathlib (the universal quantification over $N$ and
over the update sequence is kernel-checked); the proof script forms part of the paper's
verification materials.}
\end{proof}

The Jeffrey-specific content of the paper is therefore confined to the \emph{order effect}
itself (Proposition~\ref{prop:A2}), a property of probability kinematics; the first-order
protection established below rests only on separability, which any attribute-local
multiplicative update possesses.

\begin{proposition}[C1: belief-level decoupling]\label{prop:C1}
The believed-association distortion is second order---its $\bigO(c)$ coefficient is identically
zero:
\[
  \assoc(\Pbar)-\assoc(\PB)=\bigO(c^{2}).
\]
The aggregate \emph{marginal} (base-rate) distortion is first order, with
\[
  \big[\Pbar(A{=}1)-\PB(A{=}1)\big]
   = c\cdot\frac{q_0(1-q_0)\,(r_0-\beta)}{2\,\alpha\beta(1-\alpha)(1-\beta)}+\bigO(c^2),
\]
so its sign equals $\sgn(r_0-\beta)$: sign-indefinite across cue--base-rate configurations.
\anchor{Prop.~C1}
\end{proposition}

\begin{proof}
\emph{Association: the $\bigO(c)$ terms coincide by separable covariance.} For any rank-one
(separable) reweighting $D'_{ij}=g_i h_j D_{ij}/Z$ with $Z=\sum_{ij}g_i h_j D_{ij}$,
\[
  \assoc(D')=\frac{g_0 g_1 h_0 h_1}{Z^{2}}\,\assoc(D),
\]
since the cross terms in $D'_{00}D'_{11}-D'_{01}D'_{10}$ carry the identical factor
$g_0g_1h_0h_1$. By Lemma~\ref{lem:sep}, every Jeffrey route composes to exactly such a
separable reweighting (at $N=2$, a rank-one form $g_ih_j$), and the Bayes update is one
directly ($g_i=q_i/\alpha_i,\ h_j=r_j/\beta_j$). Every route therefore
maps the prior association $\assoc(P)=c$ to a \emph{positive multiple} of $c$; none can shift it
additively. Evaluating the multiplier at $c=0$---where all intermediate marginals are
$(\alpha,1-\alpha)$ and $(\beta,1-\beta)$ regardless of route---gives the same leading factor
\[
  \assoc(\PJ_{AB})=\assoc(\PJ_{BA})=\assoc(\PB)=K\,c+\bigO(c^2),
  \qquad K=\frac{q_0 q_1\, r_0 r_1}{\alpha(1-\alpha)\beta(1-\beta)} .
\]
Hence $\assoc(\Pbar)=\tfrac12\big(\assoc(\PJ_{AB})+\assoc(\PJ_{BA})\big)=Kc+\bigO(c^2)$ and
$\assoc(\Pbar)-\assoc(\PB)=\bigO(c^2)$: the $\bigO(c)$ coefficients cancel not by arithmetic
accident but because Jeffrey and Bayes are both separable, so both can only rescale the prior
association by the same leading factor.

\emph{Marginal.} Differentiating $\Pbar(A{=}1)-\PB(A{=}1)$ at $c=0$ gives the stated
coefficient; the factor $(r_0-\beta)$ governs its sign while the remaining factors are positive
on the open cube, so the sign is $\sgn(r_0-\beta)$: sign-indefinite across cue--base-rate
configurations.
\end{proof}

\begin{remark}[the mechanism]\label{rem:mechanism}
A separable reweighting can only scale an association, never manufacture one: it multiplies the
prior association by a factor, and scaling zero leaves zero. Both formalisms apply the same
scaling at leading order, so the difference enters at $\bigO(c^{2})$. The suppression is
structural, not statistical---the reason is separability, not two first-order terms cancelling
by arithmetic accident. And the structure is not Jeffrey's: by Lemma~\ref{lem:sep} the
protection would hold verbatim for any updating rule whose steps reweight through a single
attribute's partition. What is kinematic is only that the two orders \emph{differ} at all
(Proposition~\ref{prop:A2}); that their difference leaves no first-order trace in the
association is a property of attribute-locality. \anchor{Prop.~C1}
\end{remark}

\begin{remark}[economic reading]\label{rem:economic}
Consider a population of panels evaluating a common pool, each receiving the same two
impressions---one bearing on competence, one on trustworthiness---in an order idiosyncratic to
the panel. Each panel's posterior differs from the benchmark at $\bigO(c)$
(Proposition~\ref{prop:A2}). Yet the average believed association---whether the evaluators
collectively take the two attributes to go together more or less than the benchmark
warrants---carries no first-order trace of the ordering (Proposition~\ref{prop:C1}): reading
order cannot, to first order, build or erode the perceived relationship between the attributes,
because both updates reweight separably, and a separable reweighting only rescales the
association already present. What ordering shifts at first order is the level: the pool's
average perceived competence drifts at $\bigO(c)$, with sign $\sgn(r_0-\beta)$---sign-indefinite
across cue--base-rate configurations, and outside the protected class
(Proposition~\ref{prop:C3}). Picture, then, what an outside observer would find who surveyed
these panels' joint beliefs and compared them against the benchmark. Along the association
axis, the two would agree to first order: the average believed relationship between competence
and trustworthiness would sit at essentially its benchmark value, with any discrepancy of order
$c^{2}$. Along the base-rate axis they would not: the average perceived level of either
attribute would be displaced from the benchmark already at order $c$. The two coordinates of
the belief thus behave differently under the same sequencing noise---the relationship is held,
the levels are moved---and it is only in the levels that the observer could read off any trace
of the ordering at all. The footprint of sequencing lives in levels, not in the believed
relationship.
\end{remark}

\subsection{The protected class}
\label{sec:protected}

The economic reading exhibited one protected statistic and one unprotected one. Protected: the
mean believed association, whose distortion is second order. \emph{Protection}, throughout,
means exactly this---the aggregate's first-order coefficient vanishes identically, so the
statistic sits within $\bigO(c^{2})$ of its benchmark value even though every member of the
population is $\bigO(c)$ away (Propositions~\ref{prop:A2} and~\ref{prop:C1}). The association
is protected because every updating route in the model is a separable reweighting, and a
separable reweighting can only rescale an association, never shift it
(Lemma~\ref{lem:sep}, Remark~\ref{rem:mechanism}). Unprotected: the mean perceived base rate,
displaced already at $\bigO(c)$ with sign $\sgn(r_{0}-\beta)$
(Proposition~\ref{prop:C1})---separability constrains the association coordinate and leaves the
levels free to move. The split raises the question whether it is a coincidence of these two
particular summaries or a structural feature of aggregation. The answer, established next, is
the latter, and it is exhaustive: to first order, the association---and, at independence, any
smooth function of it---is the \emph{only} statistic aggregation protects; every other smooth
statistic is generically displaced at first order, exactly as the base rate was
(Proposition~\ref{prop:C3}). The base-rate drift seen above is therefore not an exception to
protection but an instance of the generic behaviour outside the protected class, and the
implication is operational: whether an aggregate registers the phenomenon at first or second
order is decided by which statistic it computes, not by anything about the population behind
it.

\begin{proposition}[C3: uniqueness of the protected statistic]\label{prop:C3}
Fix $(q_0,r_0)\in(0,1)^2$ and let $F$ be $C^1$ in a neighbourhood of $q\otimes r$ in the
probability simplex. For $\lambda\in[0,1]$ let
$\Pbar_\lambda=\lambda\,\PJ_{AB}+(1-\lambda)\,\PJ_{BA}$ be the mean belief of a population in
which a fraction $\lambda$ meets the cues in the order $A$-then-$B$.
\begin{enumerate}
\item[(i)] \textup{(Protection.)} If the differential of $F$ at $q\otimes r$, restricted to the
tangent space of the simplex, is a scalar multiple of $d\,\assoc$, then
$F(\Pbar_\lambda)-F(\PB)=\bigO(c^2)$ for every prior $(\alpha,\beta)$ and every $\lambda$.
\item[(ii)] \textup{(Uniqueness.)} Conversely, fix $\lambda\in(0,1)$. If
$F(\Pbar_\lambda)-F(\PB)=o(c)$ for all $(\alpha,\beta)$ in some open set, then $dF$ at
$q\otimes r$, restricted to the simplex, is a multiple of $d\,\assoc$.
\end{enumerate}
\anchor{Prop.~C3}
\end{proposition}

The proof is in \ref{app:proofs}; its engine is route purity: each order's leading gap is a pure
multiple of a single rank-one direction ($\kappa R_1$ or $\kappa' R_2$, inherited from
Proposition~\ref{prop:A2} and its mirror), the annihilator of the two directions is
two-dimensional and spanned by the constants and $\nabla\assoc$, and the converse follows from
the nonconstancy of $\kappa/\kappa'$ on any open set of priors \anchor{Prop.~C3}.

\begin{remark}[the protected class]\label{rem:protected}
The statistics protected to second order are exactly the \emph{local dependence measures}:
those that, to first order at independence, depend on the posterior only through its
association. This includes the association itself (Proposition~\ref{prop:C1} is the case
$F=\assoc$), the odds ratio, its logarithm, and Yule's $Q$; it excludes every marginal (the
base-rate drift of Proposition~\ref{prop:C1}), every expectation $\langle v,P\rangle$ with
nonconstant $v$ (Lemma~\ref{lem:B1}), and the headcount (Proposition~\ref{prop:C2}). Together
with Theorem~\ref{thm:B2} and Proposition~\ref{prop:C2}---which settle the non-smooth threshold
statistics, protected precisely when surplus-weighted---this yields a complete answer to
\emph{which} aggregates are second order: the dependence measures and the surplus-weighted
welfare loss, and nothing else. The mixture parameter $\lambda$ shows the statement is not an
artefact of the symmetric random-sequence population: it holds for every composition of
encounter orders. \anchor{Prop.~C3}
\end{remark}


\subsection{Incidence versus intensity}
\label{sec:criterion}

\begin{proposition}[C2: criterion split]\label{prop:C2}
Under the foregone-surplus criterion (cost $|u|$) the aggregate loss is $\bigO(c^{2})$
(Theorem~\ref{thm:B2}). If moreover the $u$-marginal is bounded and continuous at $u=0$ with
$f(0)>0$ (the same density hypothesis as Theorem~\ref{thm:B2}), then under a $0$--$1$
criterion---the \emph{incidence} of sequence-affected individuals---the aggregate is genuinely
$\bigO(c)$, not merely $o(1)$: the band mass itself scales as $f(0)\cdot\bigO(c)$. \anchor{Prop.~C2}
\end{proposition}

\begin{proof}
Throughout, fix the density hypothesis of the second part of Theorem~\ref{thm:B2}: the
$u$-marginal bounded and continuous at $u=0$ with $f(0)>0$. An individual is
sequence-affected exactly when the Jeffrey and benchmark actions differ, and by
Lemma~\ref{lem:B1} the Jeffrey score is the benchmark score perturbed by
$c\delta_{\sigma}+\bigO(c^{2})$, so the actions differ exactly when that perturbation flips
the sign of the surplus: the flip set is confined to a band of width $\bigO(c)$ at the
threshold, $\{|u|\le|c\,\delta_{\sigma}|+\bigO(c^{2})\}$, whose mass is
$f(0)\cdot\bigO(c)$---genuinely first order, not merely bounded by it, since $f(0)>0$ and
$\delta_{\sigma}$ is generically nonzero (Lemma~\ref{lem:B1}). Both criteria now integrate
over this same set and differ only in the weight they attach to its members. The incidence
weights each member by $1$ and therefore returns the band's mass: $\bigO(c)$. The surplus
criterion weights each member by its stake $|u|$, and on the band
$|u|\le|c\,\delta_{\sigma}|=\bigO(c)$ by construction: an $\bigO(c)$ mass of individuals,
each costing $\bigO(c)$, gives $\bigO(c)\cdot\bigO(c)=\bigO(c^{2})$
(Theorem~\ref{thm:B2}). The one-order gap is thus criterion arithmetic on a common set: the
extra factor of $c$ is exactly the weight $|u|$, which the band forces to first order.
\end{proof}

This is the crux for interpretation: a first-order \emph{number} of individuals are
genuinely affected, each near-marginal. The second-order quantities are the surplus-weighted
welfare loss and the believed association---the result must never be phrased as ``few are
affected,'' because the incidence is first order \anchor{Prop.~C2}.

\begin{remark}[what an audit measures]\label{rem:audit}
The two questions ``how many are affected?'' and ``what do the aggregate dependence and welfare
statistics show?'' decouple by one order in $c$. The criterion, not the phenomenon,
selects which answer an aggregate instrument returns
(Propositions~\ref{prop:C2} and~\ref{prop:C3}).
\end{remark}

\section{Scope and robustness}
\label{sec:scope}

The domain fixed in Section~\ref{sec:assumptions} by Assumptions~\ref{as:localc}--\ref{as:surplus}
is deliberately narrow, and its boundary is part of the content: each assumption is load-bearing,
and dropping it reopens a specific failure mode, located exactly below.

The binary two-attribute structure is not an analytical convenience but the canonical object of
the evaluation setting the paper addresses: a threshold decision over latent present-or-absent
attributes is the architecture of the statistical-discrimination models this framework engages
\pcite{Phelps 1972; Arrow 1973; Coate--Loury 1993}.

Attribute-locality (Assumption~\ref{as:local}) is load-bearing rather than decorative: it is exactly what makes every
Jeffrey update a separable reweighting (Lemma~\ref{lem:sep})---the engine of
Proposition~\ref{prop:C1}. A gestalt cue
bearing jointly on both attributes is out of scope, and, being non-separable, could itself
manufacture association from independence, so the immunity of Proposition~\ref{prop:A1} would
not apply.

Stated affirmatively, three things would break the results, and each is located exactly. A
non-surplus criterion breaks the welfare decoupling (Assumption~\ref{as:surplus}): under a fixed
per-error cost the aggregate loss \emph{is} the incidence and is first order
(Proposition~\ref{prop:C2})---the criterion-dependence is the content of the paper, not a
weakness of it, and Proposition~\ref{prop:C3} delimits precisely which criteria are protected. A
prior covariance outside the leading-order regime voids the separation
(Assumption~\ref{as:localc}): the results are local in $c$, and no claim is made globally. The
regime is a posture about the evaluator pool's prior, and it separates aggregation scenarios
cleanly. Where the two attributes are believed close to independent---the prior association weak
relative to its feasibility bound---the leading term is the story and the separation of scales
is informative. Where the association is entrenched---$c$ a sizable fraction of its bound, the
attributes believed to travel together almost always---the expansion's ordering of magnitudes
carries no quantitative force, and the theorems are silent: they neither protect the aggregates
there nor expose them, and silence, not failure, is the correct reading of that boundary. The
aggregation is local in a second sense as well: the population varies the encounter order (and,
for the welfare law, the stakes $u$) at a common prior, so a pool dispersed in the prior
covariance itself lies outside the stated results. And a
single decisive cue gives no effect at all (Assumption~\ref{as:soft}): the phenomenon requires
the bundle (Proposition~\ref{prop:A1}).

\section{Concluding remarks}
\label{sec:conclusion}

The paper's core insight is a separation of scales: micro-transparency does not imply
macro-detectability. An aggregate matching the Bayesian benchmark does not license the inference
that its members are Bayesian---at leading order, the population's mean believed association and
its surplus-weighted welfare loss sit within $\bigO(c^{2})$ of the benchmark while every member
diverges at $\bigO(c)$.

The instrumental reading follows from the characterization, not from any feature of data
collection. Detection turns on the criterion an instrument aggregates relative to the
phenomenon's scale, not on sample size: disaggregating a protected aggregate does not change its
order, and only a change of criterion---from the surplus-weighted loss or a dependence measure
to a statistic outside the protected class, such as the incidence or an individual-level
quantity---does. A null reading on an aggregate-built instrument is therefore not evidence that
the underlying phenomenon is null; the instrument has one order less resolution than the thing
it would detect.

Three interpretive limits should be stated once and plainly. The model exhibits a configuration
in which coherent updating produces sequence-dependent decisions; it asserts nothing about how
often that configuration is realised in actual evaluation settings. The result is
one-directional: it refutes the inference from a null aggregate to a null phenomenon, and it
validates nothing in the other direction---a null audit is uninformative both ways. And the
micro-level effect is order-idiosyncratic, reversing when the cue sequence flips, not keyed to
any attribute of the person evaluated; this is an independent reason instruments keyed to
such attributes come up empty here.

The contribution is a characterization: the class of population statistics that suppress the
first-order signature of coherent sequential updating on correlated soft cues is exactly the
local dependence measures together with the surplus-weighted welfare loss, and nothing else.
One observational consequence connects to the belief-aggregation criterion of
Section~\ref{sec:literature}: at leading order, a population's aggregate beliefs can look
externally Bayesian---order-free, benchmark-conforming on the protected statistics---even when
no member's belief is.

Two directions remain open. The suppression mechanism itself is not two-cue-specific: the
separability of attribute-local composites (Lemma~\ref{lem:sep}) holds for any number of cues,
and indeed for any attribute-local multiplicative updating rule, not Jeffrey conditioning in
particular, so the first-order protection of the dependence measures is general-$N$. What is
proved only at $N=2$ is the \emph{exact} delimitation of the protected class
(Proposition~\ref{prop:C3}); symbolic computation with general marginals verifies that it
extends without fracture at three and four cues---individual divergence first order, the
dependence and surplus-weighted aggregates second order, the incidence first order, with the
higher-way central moments, objects with no two-cue analogue, themselves protected---and what
is specific to two cues is only the one-dimensional marginal-zero perturbation space. A
general-$N$ characterization theorem and a treatment of heterogeneous and non-exchangeable cue
sequences are open.

\appendix

\section{Deferred proofs}
\label{app:proofs}

\subsection{Proof of Proposition~\ref{prop:A2}}

\emph{Coincidence at $c=0$.} At $c=0$ the prior factorises, $P(i,j)=P(A{=}i)\,P(B{=}j)$. Each
Jeffrey step then resets one coordinate's marginal to its impression while leaving the other
coordinate untouched, and the matched Bayesian update does likewise; hence
$\PJ_{AB}=\PJ_{BA}=\PB=q\otimes r$ at $c=0$. Both gaps therefore vanish at $c=0$. Each posterior
is a rational function of $c$ with denominator nonvanishing in a neighbourhood of $0$, so Taylor
expansion yields $\PJ_{AB}-\PB=cD+\bigO(c^{2})$ and
$\PJ_{AB}-\PJ_{BA}=c\,\Delta_{\mathrm{seq}}+\bigO(c^{2})$; the sequence effect carries no
constant term precisely because both sequences coincide at $c=0$.

\emph{The leading coefficients are not identically zero.} With
$Z=\alpha\beta(1-\alpha)(1-\beta)>0$, differentiating at $c=0$,
\[
  D_{00}=-D_{01}=\frac{q_0 r_0\,(q_0-\alpha)(r_0-1)}{Z},
  \qquad
  D_{10}=-D_{11}=\frac{-\,r_0\,(q_0-\alpha)(q_0-1)(r_0-1)}{Z},
\]
equivalently, in the compact rank-one form of the statement,
$D=\kappa\,q\otimes(1,-1)$ with $\kappa=(\alpha-q_0)r_0(1-r_0)/Z$. The $(1,-1)$ contrast makes
the $A$-marginal preservation manifest ($D_{i0}=-D_{i1}$), and the scalar exhibits the two ways
the gap can vanish at first order: $\kappa\propto(\alpha-q_0)$, the surprise in the
$A$-impression relative to the prior, and $\kappa\propto r_0(1-r_0)$, the spread of the
$B$-impression. Thus the leading individual gap vanishes only when the $A$-impression coincides
with the prior $A$-marginal, or the impressions are degenerate ($q_0\in\{0,1\}$ or
$r_0\in\{0,1\}$). For any $q_0\neq\alpha$ in $(0,1)$ and $r_0\in(0,1)$ we have $D\not\equiv0$.
For the sequence effect,
\[
  \Delta_{\mathrm{seq},\,00}
   =\frac{q_0 r_0\big[(\alpha-\beta)+r_0(1-\alpha)-q_0(1-\beta)\big]}{Z},
\]
whose bracket vanishes only on a single hypersurface, and which is nonzero generically: at the
witness point $\alpha=\beta=\tfrac12,\ q_0=\tfrac45,\ r_0=\tfrac3{10}$,
$\Delta_{\mathrm{seq},\,00}=-\tfrac{24}{25}\neq0$. Under full symmetry $\alpha=\beta,\ q_0=r_0$
the bracket vanishes, so the diagonal components are zero; the sequence effect itself does not
vanish there but is confined to the antisymmetric off-diagonal part,
\[
  \Delta_{\mathrm{seq},\,01}=-\Delta_{\mathrm{seq},\,10}
   =\frac{q_0(1-q_0)\,(q_0-\alpha)}{\alpha^{2}(1-\alpha)^{2}},
\]
which is zero only when additionally $q_0=\alpha$ (an uninformative impression). Both
coefficients satisfy $\sum_k D_k=\sum_k\Delta_{\mathrm{seq},\,k}=0$, as a difference of
probability distributions must. \qed

\subsection{Proof of Theorem~\ref{thm:B2}}

\emph{Decoupling: $L(c)=o(c)$.} The two actions differ iff
$\sgn(u)\neq\sgn(u+c\delta_\sigma)$. The coefficient is bounded, $|\delta_\sigma|\le M<\infty$
(it is the value-projection $\langle\vv,\partial_c(\PJ_\sigma-\PB)|_{0}\rangle$ of a function
bounded on the compact parameter domain). For fixed $\delta$ the flip set in $u$ is the open
interval strictly between $0$ and $-c\delta$, so it is contained in the \emph{punctured} band
\[
  B_c:=\{\,u:\ 0<|u|\le cM\,\},
\]
on which the cost obeys $|u|\le cM$. With $\mu$ the marginal law of $u$, Tonelli gives
\[
  L(c)=\iint_{\mathrm{flip}}|u|\,g(u,\delta)\,du\,d\delta
   \;\le\; cM\!\iint_{u\in B_c}\! g(u,\delta)\,du\,d\delta
   \;=\; cM\,\mu(B_c).
\]
The bands are nested and decrease as $c\downarrow0$ to $\bigcap_{c>0}B_c=\varnothing$: any
$u\neq0$ leaves the band once $cM<|u|$, and $u=0$ is excluded from every $B_c$. As $\mu$ is a
probability measure (finite), continuity from above gives $\mu(B_c)\to\mu(\varnothing)=0$, i.e.\
$\mu(B_c)=o(1)$. Hence
\[
  \frac{L(c)}{c}\le M\,\mu(B_c)\longrightarrow0,\qquad L(c)=o(c),
\]
for \emph{any} integrable density, and irrespective of an atom at the threshold---the atom sits
at the excised point $u=0$, where the cost $|u|$ vanishes, so it never enters $B_c$.

\emph{Named constant under continuity.} If in addition the $u$-marginal is continuous at $0$
with $f(0)>0$, the flip-set integral for fixed $\delta$ is
$\int_0^{|c\delta|}t\,g(0,\delta)\,dt=g(0,\delta)\,\delta^{2}c^{2}/2+\bigO(c^3)$; integrating
over $\delta$ with $f(0)=\int g(0,\delta)\,d\delta$ and
$\E[\delta^2\mid u=0]=\int\delta^2 g(0,\delta)\,d\delta/f(0)$ gives
$L(c)=\tfrac{f(0)}{2}\E[\delta_\sigma^2\mid u=0]\,c^{2}+\bigO(c^3)$. \qed

\subsection{Proof of Proposition~\ref{prop:C3}}

\emph{Route purity.} By Proposition~\ref{prop:A2}, $\PJ_{AB}-\PB=c\,\kappa R_1+\bigO(c^2)$ with
$R_1=q\otimes(1,-1)$ and $\kappa=(\alpha-q_0)r_0(1-r_0)/Z$; exchanging the roles of the cues,
$\PJ_{BA}-\PB=c\,\kappa' R_2+\bigO(c^2)$ with $R_2=(1,-1)\otimes r$ and
$\kappa'=(\beta-r_0)q_0(1-q_0)/Z$. Hence $\Pbar_\lambda-\PB=c\,M_\lambda+\bigO(c^2)$ with
$M_\lambda=\lambda\kappa R_1+(1-\lambda)\kappa' R_2$, and by the chain rule
\[
  F(\Pbar_\lambda)-F(\PB)\;=\;c\,\bigl\langle \nabla F(q\otimes r),\,M_\lambda\bigr\rangle+\bigO(c^2).
\]

\emph{The annihilator.} The system $\langle V,R_1\rangle=\langle V,R_2\rangle=0$ in
$V\in\mathbb{R}^{2\times2}$ has a two-dimensional solution space. The constant matrix $J$ (all
entries $1$) solves it because the rows of $R_1$ and the columns of $R_2$ sum to zero; the
gradient of the association at the independent posterior,
\[
  \nabla\assoc(q\otimes r)=
  \begin{pmatrix} q_1r_1 & -q_1r_0\\ -q_0r_1 & q_0r_0\end{pmatrix},
\]
solves it by direct computation
($\langle\nabla\assoc,R_1\rangle=q_0q_1r_1-q_0q_1r_0-q_1q_0r_1+q_1q_0r_0=0$, and symmetrically
for $R_2$). Since $J$ and $\nabla\assoc(q\otimes r)$ are linearly independent, they span the
annihilator.

\emph{(i)} If $dF|_{q\otimes r}$ restricted to the simplex is a multiple of $d\,\assoc$, then
$\nabla F(q\otimes r)=xJ+y\,\nabla\assoc(q\otimes r)$ for scalars $x,y$, so
$\langle\nabla F,M_\lambda\rangle=0$ identically in $(\alpha,\beta,\lambda)$.

\emph{(ii)} Write $x=\langle\nabla F,R_1\rangle$ and $y=\langle\nabla F,R_2\rangle$, so the
first-order coefficient is $\lambda\kappa\,x+(1-\lambda)\kappa'\,y$, vanishing on an open set
$U$ of priors. If $y\neq0$, then on $U\setminus\{\beta=r_0\}$ the ratio $\kappa/\kappa'$ would
equal the constant $-(1-\lambda)y/(\lambda x)$ (with $x\neq0$; if $x=0$ then $\kappa'y\equiv0$
forces $y=0$ at once, since $\kappa'$ vanishes only on $\{\beta=r_0\}$). But
$\kappa/\kappa'=\dfrac{(\alpha-q_0)\,r_0(1-r_0)}{(\beta-r_0)\,q_0(1-q_0)}$ has nonvanishing
$\alpha$-derivative wherever $\beta\neq r_0$, so it is nonconstant on any open set: a
contradiction. Hence $y=0$, and then $\kappa x\equiv0$ on $U$ forces $x=0$. So
$\nabla F(q\otimes r)$ lies in the annihilator $\mathrm{span}\{J,\nabla\assoc\}$; restricted to
the tangent space of the simplex, where $J$ acts as zero, $dF$ is a multiple of $d\,\assoc$.
\qed

\section*{Declaration of Generative AI and AI-Assisted Technologies in the Manuscript Preparation Process}

%% DRAFT — Anurag to review, edit, and take responsibility per Elsevier JME author guide.
During the preparation of this work the author used Anthropic's Claude (large language model)
in order to: draft and refine manuscript prose; propose, prove, and cross-verify
Proposition~\ref{prop:C3} (uniqueness of the protected statistic), which was then independently
re-derived and machine-verified by the author; propose bibliographic sources for
positioning against neighbouring literatures; and structure the manuscript to conform to the
target journal's house style. After using this tool, the author reviewed and edited the
content as needed, independently verified the mathematical statements, and takes full
responsibility for the content of the published article.

\section*{Declaration of competing interest}

%% DRAFT — Anurag to confirm wording against the current Elsevier/JME author guide at submission.
The author declares that he has no known competing financial interests or personal relationships
that could have appeared to influence the work reported in this paper.

\section*{Acknowledgements}

%% To be completed by author.

\bibliographystyle{elsarticle-harv}
\bibliography{bibliography}

\end{document}
